% root file is THHBP2.tex

\section{Introduction}
Topological Hochschild homology has applications to many areas of mathematics, for example deformations of $\mathbb{A}_{\infty}$-algebras~\cite{Laz01}, string topology~\cite{CJ02}, algebraic K-theory~\cite{BHM93}, $p$-adic Hodge theory~\cite{BMS19}  and knot Floer homology~\cite{LOT15}. It also played an important role in the recent disproof of the telescope conjecture~\cite{BHLS23}. The topological Hochschild homology of $\mathbb{F}_p$ can be used to provide a new proof of Bott periodicity~\cite{HN20} and variants on topological Hochschild homology have been used to provide a new proof of the Segal conjecture for cyclic groups of prime order~\cite{HW21,AKZ25}. 
We therefore consider the topological Hochschild homology of ring spectra to be of fundamental importance. 

In homotopy theory, in addition to characteristic zero and characteristic $p$, there are characteristics corresponding to pairs $(n,p)$ where $n$ is called the height and $p$ is a prime. This also leads to higher height analogues of rings of integers in local fields of which truncated Brown--Peterson spectra are an example~\cite{Rog25}. The topological Hochschild homology of the integers $\mathbb{Z}$, and more generally rings of integers in number fields, has been related to special values of Dedekind zeta functions in work of Morin~\cite{Mor24}, which builds on work of Bhatt--Morrow--Scholze~\cite{BMS19}. To generalize this to characteristics corresponding to higher heights $n\ge 1$, a complete understanding of topological Hochschild homology of truncated Brown--Peterson spectra $\BP\langle n\rangle$ is desirable.



When $n=-1$ then $\BP\langle -1\rangle=\mathbb{F}_p$ and when $n=0$ then $\BP\langle 0\rangle=\mathbb{Z}_{(p)}$ and their topological Hochschild homology was computed by B\"okstedt~\cite{Bok85}. When $n=1$, then $\BP\langle 1\rangle=\ell$ is the Adams summand of topological complex K-theory $\mathrm{ku}$, which is built from the classifying space of the infinite unitary group, which represents complex vector bundles. In this case, $\THH_*(\ell)$ was completely computed by~\cite{AHL10} and $\THH_*(\ku)$ was computed by~\cite{Lee22} and \cite{Cha25} using different methods. Beyond these examples, there are no known complete computations of $\THH_*(R)$ where $R$ is a ring spectrum that is not Eilenberg--MacLane.

To compute topological Hochschild homology of truncated Brown--Peterson spectra, the traditional approach is to compute the topological Hochschild homology with coefficients in a finite field $\mathbb{F}_p$ and then use a lattice of Bockstein spectral sequences to reduce down to topological Hochschild homology of truncated Brown--Peterson spectra. Here one uses the fact that the coefficients of truncated Brown--Peterson spectra are $\BP\langle n\rangle_*=\mathbb{Z}_{(p)}[v_1,\cdots ,v_n]$ and we can therefore consider $v_i$-Bockstein spectral sequences such as 
\[ 
\THH_*(\BP\langle n\rangle;\BP\langle n\rangle/v_i)[v_i]\implies \THH_*(\BP\langle n\rangle)\,.
\]

The first step is therefore the computation 
\[ \THH_*(\BP\langle n\rangle;\mathbb{F}_p)=\mathbb{F}_p[\mu^{p^{n+1}}]\langle \lambda_1,\dots ,\lambda_{n+1} \rangle
\]
first appearing in~\cite{ACH24} (cf.~\cite{Ang26}). Here  $\mathbb{F}_p\langle \lambda_1,\cdots ,\lambda_n\rangle$ denotes an exterior algebra on generators in degree $|\lambda_{i}|=2p^i-1$ for $1\le i\le n+1$ and $\mathbb{F}_p[\mu^{p^{n+1}}]$ is a polynomial algebra with generator in degree $|\mu^{p^{n+1}}|=2p^{n+1}$. From there the first author computed 
\[ \THH_*(\BP\langle n\rangle ; \mathbb{Z}_{(p)})_p^{\wedge}\]
with Culver and H\"oning for arbitrary heights $n$ as well as 
\[ \THH_*(\BP\langle 2\rangle ; M)\]
when $M\in \{\BP\langle 2\rangle/(p,v_1),\BP\langle 2\rangle /(p,v_2)\}$. The next goal in this program is to compute 
\[ \THH_*(\BP\langle 2\rangle;\BP\langle 1\rangle) \,,
\]
which we accomplish in this paper. Throughout, we write $\BP\langle n\rangle$ for a $\mathbb{E}_3$-$\MU$-algebra form of truncated Brown--Peterson spectra at the prime $p$. 
\begin{thmx}\label{thm:main}
  The topological Hochschild homology of $\BP\langle 2\rangle$ is a direct sum
  \begin{equation}
    \THH_*(\BP \langle 2 \rangle ; \BP \langle 1 \rangle) \cong \mathcal{F}\oplus \Sigma^{2p^2-1}\mathcal{F} \oplus \mathcal{T} \oplus \Sigma^{2p - 1} \mathcal{T}
  \end{equation}
  where $\mathcal{T}$ is the $\Z_{(p)}[v_1]$-submodule of $\THH_*(\BP \langle 1 \rangle)$ containing the torsion elements that are in the image of the multiplication by $v_1^p$ map and $\mathcal{F}$ denotes the torsion free part of $\THH_*(\ell)$. 
\end{thmx}

As part of this program, we present a new computational tool, which is in some sense an instance of the Brun spectral sequence considered by~\cites{Bru00,Hon20}. This computational tool gives an inductive approach to computing $\THH(\BP\langle n\rangle)$. First one computes $\THH(\BP\langle n-1\rangle)$ along with the cap product of a certain class in the of topological Hochschild cohomology $\BP\langle n-1\rangle$. Then one uses the Brun spectral sequence from this paper to compute $\THH(\BP\langle n\rangle;\BP\langle n-1\rangle)$.  Finally, one computes the remaining $v_{n}$-Bockstein spectral sequence to compute $\THH_*(\BP\langle n\rangle)$, which we leave to future work. 

As an application of our work, we show that $\BP\langle n\rangle$ is not a Thom spectrum of a $2$-fold loop map over the sphere spectrum for $n\ge 2$ at the prime $p=2$. 

\begin{thmx}\label{thm:main-thom}
Truncated Brown--Peterson spectra $\BP\langle n\rangle$ are not Thom spectra of $2$-fold loop maps for any $n\ge 2$ at the prime $p=2$. 
\end{thmx}

\subsection{Outline}
We first recall the necessary results from~\cite{AHL10}, \cite{AR05}, and \cite{ACH24} that we build on in this paper in Section~\ref{recollections}. 


\subsection{Conventions}
We write $L_{E}$ for Bousfield localization at a spectrum $E$. We write $\BP\langle n\rangle$ for a family of $\bE_3$-$\MU$-algebra forms of $\BP\langle n\rangle$ such that there are maps
\[
\BP\to \MU \to  \dots \to \BP\langle n\rangle\to  \BP\langle n-1\rangle\to  \dots  \to \bZ_{(p)}\to \bF_{p}
\]
of $\bE_3$-rings. 
We fix classes $v_i$ such that on graded commutative rings
\[
\BP_*\to  \BP\langle n\rangle_*
\]
is given by sending $v_i$ to $v_i$ for $0\le i\le n$ (with $v_0=p$) and $v_j\mapsto 0$ otherwise. 
This also fixes the map of graded commutative rings 
\[  
    \BP\langle n\rangle_*\to  \BP\langle n-1\rangle_*
\]
sending $v_i$ to $v_i$ for $0\le i\le n-1$ and $v_n\mapsto 0$. Such a family exists by~\cite{HW22}*{Theorem~A}. 
Alternatively, when $n\le 2$ and $p\in \{2,3\}$, there are $\bE_2$-$\MU$-algebra forms of $\BP\langle 2\rangle$ denoted $\tmf_1(3)$ and $\taf^{D}$ respectively, which have the extra property that their $\bE_3$-ring structures lift to $\bE_\infty$-ring structures by~\cites{HL10,LN12,LN14,CM15}. 

Recall that if $R$ is an $\mathbb{E}_1$-ring and $M$ is an $R\otimes R^{\op}$-module, then we can define 
\[ 
\THH(R;M):=M\otimes_{R\otimes R^{\op}}R
\]
and wen $M=R$ we simply write $\THH(R):=\THH(R;R)$. Here we write $\otimes$ for the derived smash product in spectra. We also write $\otimes$ for the tensor product of abelian groups and we let context dictate the meaning. We also define 
\[ 
\THC(R;M):=\map_{R\otimes R^{\op}}(M,R)
\]
where we write $\map_{R\otimes R^{\op}}$ for the mapping spectrum in $R\otimes R^{\op}$-modules. 

We write $\dot{=}$ for an equality that holds up to multiplying one side by a unit. 

\subsection{Acknowledgements}
The first author would like to thank Dominic Leon Culver and Eva H\"oning for their contributions during an earlier collaboration. The first author would also like to thank Haldun \"Ozg\"ur Bayindir for helpful conversations related to this project. The authors would also like to thank the organizers of the the IWoAT workshop in 2025 in Hangzhou, China where the authors began this collaboration. 